% !TEX root = ../main.tex
\chapter{変数・引数・任意の値}
\label{chap:ch07_variables_and_forall}

\begin{chapterabstract}
本章では、値を一つに固定しない主張を扱う。
「任意の価格について」「すべてのユーザーについて」という形の仕様を書くために、
引数を取る定理、\texttt{forall}、証明を始める \texttt{intro} を導入する。
あわせて、無名関数・型パラメータ・関数合成と、\texttt{Option} と \texttt{match} を使った小さなパイプラインも扱い、
後続の圏論の章への橋渡しとする。
\end{chapterabstract}

\begin{learninggoals}
\item 引数を取る定理で「任意の入力について成り立つ」主張を書ける。
\item \texttt{forall} と \texttt{intro} を使い、すべての値についての仕様を証明できる。
\item 無名関数 \texttt{fun}、型パラメータ、関数合成を読める。
\item \texttt{Option} と \texttt{match} を使った小さなパイプラインを読める。
\end{learninggoals}

\section{この章を読む理由}

テストは、選んだ具体例についての確認である。
一方、仕様で本当に言いたいのは、しばしば「すべての入力でこの性質が成り立つ」という主張である。
本章では、値を固定しない主張の書き方と証明の始め方を学ぶ。
これは、後のマイグレーション証明や圏論の法則確認で何度も使う基本形になる。

\section{任意の値についての定理}

定理が引数を取ると、その引数は「任意の値」を表す。
次は、税額が \texttt{0} のとき、税込金額が元の価格と一致することを、任意の \texttt{price} について示す定理である。

\begin{listing}[htbp]
\begin{minted}{lean4}
def addTax (price tax : Nat) : Nat :=
  price + tax

theorem addTax_zero (price : Nat) :
    addTax price 0 = price := by
  rfl
\end{minted}
\caption{\texttt{addTax\_zero}}
\label{lst:ch07_forall_price}
\end{listing}

コード\ref{lst:ch07_forall_price}の \texttt{addTax\_zero} は、特定の価格ではなく、任意の \texttt{price} についての主張である。
\texttt{addTax price 0} は定義により \texttt{price + 0} になり、これは \texttt{price} に展開されるので、\texttt{rfl} で閉じられる。

\section{\texttt{forall} と \texttt{intro}}

「すべての \texttt{x} について \texttt{P x}」という主張は、\texttt{forall x, P x} と書ける。
これを証明するには、任意の値を一つ \texttt{intro} で受け取り、その値について示す。
次は、ID を取り出す処理が値を変えない、という仕様である。

\begin{listing}[htbp]
\begin{minted}{lean4}
def toDtoId (id : Nat) : Nat :=
  id

def PreservesId (f : Nat -> Nat) : Prop :=
  forall id, f id = id

theorem toDtoId_preservesId : PreservesId toDtoId := by
  intro id
  rfl
\end{minted}
\caption{\texttt{PreservesId}}
\label{lst:ch07_forall_intro}
\end{listing}

コード\ref{lst:ch07_forall_intro}の \texttt{PreservesId} は「すべての \texttt{id} について \texttt{f id = id}」という仕様である。
\texttt{intro id} で任意の ID を一つ受け取り、\texttt{toDtoId} の定義から値がそのまま返るので、\texttt{rfl} で閉じる。

\section{高階関数と関数合成を段階的に読む}

\subsection{無名関数と、関数を受け取る関数}

関数はその場で名前なしに書ける。
\texttt{fun n => n + 1} は「\texttt{n} を受け取り \texttt{n + 1} を返す関数」である。
まずは\texttt{inc}と\texttt{applyTwice}だけを読み、値の引数と関数の引数を区別する。

\subsection{型パラメータと関数合成}

型を引数に取る関数合成は、出力の型と入力の型が合う関数をつなぐ。

\begin{listing}[htbp]
\begin{minted}{lean4}
def inc : Nat -> Nat :=
  fun n => n + 1

def applyTwice (f : Nat -> Nat) (x : Nat) : Nat :=
  f (f x)

def compose {A B C : Type} (g : B -> C) (f : A -> B) : A -> C :=
  fun x => g (f x)

#eval applyTwice inc 10
#eval compose inc inc 10
\end{minted}
\caption{\texttt{inc}・\texttt{applyTwice}・\texttt{compose}}
\label{lst:ch07_compose}
\end{listing}

コード\ref{lst:ch07_compose}では、\texttt{inc} は無名関数に名前を付けた定義、\texttt{applyTwice} は関数を引数に取る関数である。
\texttt{compose} は \texttt{\{A B C : Type\}} という型パラメータを取り、\texttt{compose g f} を「先に \texttt{f}、後で \texttt{g}」と読む。
波括弧の型パラメータは、呼び出し側の型から Lean が推論する。

\begin{leanbox}
関数合成で大切なのは、型がつながることである。
\texttt{A -> B} の出力 \texttt{B} が \texttt{B -> C} の入力と一致するからつなげられる。
この見方が、後の章で扱う「射の合成」につながる。
\end{leanbox}

\begin{exercisebox}
\texttt{compose g f}の途中の型を書き、なぜ\texttt{g}を先に適用できないかを説明する。答えでは、\texttt{f : A -> B}と\texttt{g : B -> C}の接続点\texttt{B}を明示する。
\end{exercisebox}

図\ref{fig:ch07-pipeline}では、各関数の出力型を次の関数の入力型へつなぎ、複数の文字列変換を一つの処理として読む。

\FigurePlaceholder{fig-ch07-pipeline.png}{0.90}{文字列変換の合成}{fig:ch07-pipeline}

\section{\texttt{Option} と \texttt{match} の小さなパイプライン}

値が「あるかもしれない／ないかもしれない」ことは \texttt{Option} で表す。
値がある場合は \texttt{some x}、ない場合は \texttt{none} と書き、\texttt{match} で場合分けする。

\begin{listing}[htbp]
\begin{minted}{lean4}
def surround (s : String) : String :=
  "[" ++ s ++ "]"

def nonempty (s : String) : Option String :=
  if s == "" then none else some s

def decorateIfNonempty (s : String) : Option String :=
  match nonempty s with
  | none => none
  | some value => some (surround value)

def decorateAll (xs : List String) : List (Option String) :=
  xs.map decorateIfNonempty

#eval decorateIfNonempty ""
#eval decorateIfNonempty "lean"
#eval decorateAll ["a", "", "b"]
\end{minted}
\caption{\texttt{decorateIfNonempty} と \texttt{decorateAll}}
\label{lst:ch07_pipeline}
\end{listing}

コード\ref{lst:ch07_pipeline}では、\texttt{nonempty} が空文字列を \texttt{none}、それ以外を \texttt{some s} に変換する。
\texttt{decorateIfNonempty} は \texttt{match} で \texttt{Option} を場合分けし、値があるときだけ装飾する。
\texttt{decorateAll} は \texttt{xs.map} で各要素に同じ処理を適用する。

\begin{exercisebox}
\texttt{decorateIfNonempty "lean"}を、\texttt{nonempty}の結果、\texttt{match}で選ばれる枝、最終的な\texttt{Option String}の順に追跡する。\texttt{List.map}は各要素へ同じ関数を適用するだけであり、第14章で扱う関手則まで本節が証明しているわけではない。
\end{exercisebox}

\section{本章のまとめ}

\begin{summarybox}
要点と記法
\begin{itemize}
\item \texttt{forall x, P x} / \texttt{intro x} \quad 全称量化と、その導入。
\item \texttt{fun x => 本体} \quad 無名関数。
\item \texttt{\{A : Type\}} \quad 推論される型パラメータ。
\item \texttt{some} / \texttt{none} / \texttt{match} / \texttt{xs.map} \quad \texttt{Option} と一括適用。
\end{itemize}
\end{summarybox}

\section{次章への接続}

ここまでの証明は \texttt{rfl} だけで閉じてきた。
次章では、\texttt{rfl} だけでは届かない場面のために、\texttt{simp} と \texttt{rw} という次の道具を導入する。
\texttt{simp} は明らかな整理を任せ、\texttt{rw} は既知の等式で書き換える。
